\begin{tikzpicture}[x={\dimexpr\linewidth/169\relax},y=1mm,
  font=\figurefont\footnotesize,text=NNInk]
  \nnTensorState{in}{10}{0}{NNInput}{输入}
  \nnTensorState{out}{160}{0}{NNOutput}{输出}
  \node[align=center,anchor=north] at (10,-13) {$28\times28$\\$64$通道};
  \node[align=center,anchor=north] at (160,-13) {$28\times28$\\$80$通道};
  \nnConvOp{b1}{58}{27}{1\times1}
  \nnConvOp{r3}{58}{9}{1\times1}
  \nnConvOp{b3}{97}{9}{3\times3}
  \nnConvOp{r5}{58}{-9}{1\times1}
  \nnConvOp{b5}{97}{-9}{5\times5}
  \node[nn operator,minimum width=13mm,minimum height=13mm,align=center]
    (pool) at (58,-27) {Max\\$3\times3$};
  \nnConvOp{bp}{97}{-27}{1\times1}
  \nnConcatOp{cat}{135}{0}
  % 同一输入经平滑曲线送入四个分支；输出仍按通道拼接。
  \foreach \target in {b1,r3,r5,pool}{
    \draw[nn flow] (in.east) to[out=0,in=180,looseness=1.05] (\target.west);
  }
  \foreach \a/\b/\cc in {r3/b3/16,r5/b5/8,pool/bp/64}{
    \draw[nn operator path] (\a.east)--node[above=1mm] {$\cc$} (\b.west);
  }
  \draw[nn operator path] (b1.east)--node[above=1mm] {$16$} (108,27)
    .. controls (121,27) and (126,10) .. (cat.135);
  \draw[nn operator path] (b3.east)--node[above=1mm] {$32$} (113,9)
    .. controls (121,9) and (125,3.5) .. (cat.165);
  \draw[nn operator path] (b5.east)--node[above=1mm] {$16$} (113,-9)
    .. controls (121,-9) and (125,-3.5) .. (cat.195);
  \draw[nn operator path] (bp.east)--node[above=1mm] {$16$} (113,-27)
    .. controls (122,-27) and (126,-10) .. (cat.225);
  \draw[nn operator path,draw=NNOutput] (cat.east)--(out.west);
\end{tikzpicture}
